\scanpage{9}
\mspage{3}
But how to verify that $(p;g)$ define not only an endomorphism, but an automorphism? That seems delicate\ldots

In the case of $\tau$, we have $p=-1$, and $g=l_1$:
\[
\tau\rho=\operatorname{int}(l_1)\rho\tau
\]
We want to express the commutation of $u$ (as an outer automorphism \inserted{respecting the structure with loops}) not only with $\rho$ i.e. $\unsure{\mathfrak S_3^+|\mathcal S_3^+}$,
\margin{\uline{NB} This implies that it acts \unsure{trivial\textsuperscript{ly}} on the set of the three conjugacy classes of loop subgroups --- what precedes applies.}
but with $\unsure{\mathfrak S_3|\mathcal S_3}$ in its entirety. It suffices to express the commutation modulo outer automorphisms \unsure{with|:} $\sigma_\infty$, i.e. to write that there exists $h\in\pi$ (also unique) such that
\[
u\sigma_\infty=\operatorname{int}(h)\deleted{\unsure{u|\tau}}\sigma_\infty u,
\]
and it will suffice to express it for the arguments $l_0$ and $l_1$. We find the conditions
\[
\underset{\substack{\Vert\\\underbrace{u(l_1)}_{\operatorname{int}(g)l_1^p}}}{u\sigma_\infty(l_0)}=\operatorname{int}(h)\underbrace{\sigma_\infty(l_0^p)}_{l_1^p}
\qquad\text{i.e.}\quad\operatorname{int}(g^{-1}h)l_1^p=l_1^p\quad\text{i.e.}
\]
\[
g^{-1}h\in L_1\quad\text{i.e.}\quad \begin{gathered}[t]h=gl_1^q\\q\in\hat{\mathbb Z}\end{gathered}
\]
and
\[
u\underbrace{\sigma_\infty(l_1)}_{l_0}=\operatorname{int}(h)\sigma_\infty\underbrace{u(l_1)}_{\operatorname{int}(g)l_1^p}\qquad\text{i.e.}
\]
\[
l_0^p=\operatorname{int}(h\sigma_\infty(g))l_0^p\qquad\text{or}
\]
\[
l_0^p=\operatorname{int}(gl_1^q\sigma_\infty(g))l_0^p\qquad\text{which is equivalent to}
\]
\[
\deleted{\operatorname{int}(\unsure{gl_1^q}\sigma_\infty g)L_1=L_0}\qquad\text{i.e. the condition}
\]
\[
\exists_q\ \text{such that}\quad gl_1^q\sigma_\infty(g)\in L_0
\]
\mspage{8}
